\documentclass[a4paper,12pt,oneside]{report}

% ==========================================
% 1. CÁC GÓI CƠ BẢN VÀ ĐỊNH DẠNG TRANG
% ==========================================
\usepackage[utf8]{vietnam}              % Hỗ trợ gõ tiếng Việt có dấu
\usepackage{indentfirst}                % Thụt lề đoạn đầu tiên của chương/mục
\usepackage[top=2cm,inner=3cm,outer=2cm,bottom=2cm,headheight=15pt]{geometry} % Cấu hình lề
\usepackage{etoolbox}                   % Cung cấp lệnh \patchcmd để sửa cấu trúc mặc định của LaTeX

% ==========================================
% 2. XỬ LÝ HÌNH ẢNH, BẢNG BIỂU & TOÁN HỌC
% ==========================================
\usepackage{graphicx}                   % Dùng để chèn hình ảnh
\graphicspath{ {images/} }              % Thư mục mặc định chứa hình ảnh
\usepackage{float}                      % Cho phép dán cứng vị trí ảnh với thuộc tính [H]
\usepackage{tabularx}                   % Dùng cho bảng có độ rộng tùy chỉnh
\usepackage{longtable}                  % Hỗ trợ chèn bảng dài vượt quá 1 trang
\usepackage{booktabs}                   % Đường kẻ bảng chuẩn học thuật (\toprule, \midrule, \bottomrule)
\usepackage{makecell}                   % Tinh chỉnh header hoặc ô trong bảng
\usepackage{pdflscape}                  % Hỗ trợ xoay ngang trang giấy (cho bảng/hình rộng)
\usepackage{amsmath}                    % Hỗ trợ viết các công thức toán học chuyên sâu
\usepackage{tcolorbox}                  % Hỗ trợ tạo các hộp màu (dùng cho trang bìa)
\renewcommand\theadfont{\bfseries}      % Mặc định in đậm các dòng Header của bảng
\usepackage{dirtree}
\usepackage[hidelinks]{hyperref} % Ẩn hoàn toàn khung và màu, vẫn click được
\newcommand{\degree}{\ensuremath{^{\circ}}}

% ==========================================
% 3. CẤU HÌNH TÀI LIỆU THAM KHẢO & ĐƯỜNG DẪN
% ==========================================
\usepackage[square]{natbib}             % Quản lý tài liệu tham khảo
\setcitestyle{super}                    % Trích dẫn hiển thị dạng số mũ trên cao (vd: [1])
\usepackage{url}                        % Cho phép chèn link URL

% ==========================================
% 4. CHÈN CODE & CẤU TRÚC SUBFILE (FILE CON)
% ==========================================
\usepackage{subfiles}                   % Cho phép chia project thành nhiều file .tex nhỏ
\usepackage{fancyvrb}                   % Gói để hiển thị code (có format)
\fvset{
	tabsize=4,                          % Khoảng cách tab trong code
	frame=single,                       % Thêm khung hình chữ nhật bọc đoạn code
}

% ==========================================
% 5. ĐIỀU CHỈNH KHOẢNG CÁCH TIÊU ĐỀ CHƯƠNG
% ==========================================
% Xóa khoảng trống thừa (rất rộng) ở đầu trang bị mặc định bởi LaTeX cho các thẻ \chapter
\makeatletter
\def\@makechapterhead#1{%
  {\parindent \z@ \raggedright \normalfont
    \interlinepenalty\@M
    \ifnum \c@secnumdepth >\m@ne
        \Huge\bfseries \@chapapp\space \thechapter:\space #1\par\nobreak
    \else
        \Huge\bfseries #1\par\nobreak
    \fi
    \vskip 20\p@ % Khoảng cách từ \chapter xuống đoạn văn bản
  }}
\def\@makeschapterhead#1{%
  {\parindent \z@ \raggedright
    \normalfont
    \interlinepenalty\@M
    \Huge \bfseries  #1\par\nobreak
    \vskip 20\p@ % Khoảng cách từ \chapter* (như Lời Cảm Ơn) xuống văn bản
  }}
\makeatother

% ==========================================
% 6. TẠO DÒNG KẺ TRỐNG CHỖ (Cho giảng viên ghi tay)
% ==========================================
\usepackage{pgffor, ifthen}
\newcommand{\notes}[3][\empty]{%
	\vspace*{10pt}%
	\foreach \n in {1,...,#2}{%
		\noindent\makebox[#3][l]{\dotfill}\par\vspace{8pt}%
	}%
	\vspace{2cm}
	\hfill
	\begin{minipage}{0.45\textwidth}
		\centering
		\textbf{Ký tên}\\[0.5cm]
		Họ và tên
	\end{minipage}
}

% ==========================================
% 7. TÙY CHỈNH HEADER (TIÊU ĐỀ) VÀ FOOTER (CHÂN TRANG)
% ==========================================
\usepackage{fancyhdr}                   % Gói tạo header & footer tùy chỉnh
\usepackage{emptypage}                  % Ẩn số trang & header ở các trang trắng

% --- A. Style Khởi Đầu (Trang Lời Cảm Ơn, Nhận Xét...) ---
\fancypagestyle{plain}{
    \fancyhf{}                          % Xóa trắng mọi thứ
    \fancyfoot[R]{\thepage}             % Chỉ hiện số trang ở góc dưới bên Phải (R)
    \renewcommand{\headrulewidth}{0pt}  % Khong có đường gạch ngang ở Header
}

% --- B. Style Nội Dung (Từ Mục Lục Trở Về Sau) ---
\fancypagestyle{fancy}{
	\fancyhf{}                          % Phải xóa trắng trước khi thiết lập
	% --- Cấu hình Header (Trên Cùng) ---
	\lhead{\fontsize{9}{11} \selectfont \MakeUppercase{Xây dựng hệ thống thử kiểu tóc ảo sử dụng mô hình khuếch tán Stable Diffusion}}
	\rhead{}                            % Góc trên Phải: Trống
	
	% --- Cấu hình Footer (Dưới Cùng) ---
	\lfoot{\fontsize{10}{12} \selectfont Nguyễn Tiến Phát - B2306569}  % Góc dưới Trái
	\rfoot{\fontsize{10}{12} \selectfont \thepage}                     % Góc dưới Phải hiển thị số trang
	
	% --- Đường kẻ phân cách ---
	\renewcommand{\headrulewidth}{0.4pt}  % Kẻ 1 đường thẳng dưới Header
	\renewcommand{\footrulewidth}{0.4pt}  % Kẻ 1 đường thẳng trên Footer
}

% ==========================================
% 8. BẮT ĐẦU TÀI LIỆU
% ==========================================
\begin{document}

% ------------------- PHẦN BÌA -------------------
\pagenumbering{gobble}          % Tạm tắt đánh số trang cho phần bìa
\subfile{bia-ngoai}             % Gọi thiết kế bìa cứng
\subfile{bia-trong}             % Gọi thiết kế bìa trong

\clearpage                      % Sang trang mới
% ------------------- PHẦN LỜI CẢM ƠN & NHẬN XÉT ---
\pagenumbering{roman}           % Đánh số trang kiểu la mã: i, ii, iii...
\pagestyle{plain}               % Dùng style trơn (chỉ có số trang)

\chapter*{LỜI CẢM ƠN}
\phantomsection
\addcontentsline{toc}{chapter}{LỜI CẢM ƠN}
\subfile{chapters/loi-cam-on}

\chapter*{TÓM TẮT}
\phantomsection
\addcontentsline{toc}{chapter}{TÓM TẮT}
\subfile{chapters/tom-tat}

\chapter*{ABSTRACT}
\phantomsection
\addcontentsline{toc}{chapter}{ABSTRACT}
\subfile{chapters/abstract}

\chapter*{NHẬN XÉT CỦA GIẢNG VIÊN}
\phantomsection
\addcontentsline{toc}{chapter}{NHẬN XÉT CỦA GIẢNG VIÊN}
\subfile{chapters/nhan-xet-gv}

% ------------------- PHẦN MỤC LỤC -----------------
\clearpage                      
\pagestyle{fancy}               % >>> BẬT STYLE FANCY (FULL HEADER/FOOTER) TỪ ĐÂY MÃI VỀ SAU <<<

% Thuốc đặc trị: Mặc định LaTeX ép các trang Mục Lục hay đầu Chương là kiểu trơn (plain).
% Ta ghi đè style plain thành fancy từ đây trở đi để tất cả đều có Header/Footer.
\makeatletter
\let\ps@plain\ps@fancy
\makeatother

\tableofcontents                % Xuất Mục Lục chính

\clearpage
\phantomsection
\addcontentsline{toc}{chapter}{DANH MỤC HÌNH ẢNH}
\listoffigures                  % Xuất Danh mục hình ảnh

\clearpage
\phantomsection
\addcontentsline{toc}{chapter}{DANH MỤC BẢNG BIỂU}
\listoftables                   % Xuất Danh mục bảng biểu

\chapter*{DANH SÁCH TỪ VIẾT TẮT}
\phantomsection
\addcontentsline{toc}{chapter}{DANH MỤC VIẾT TẮT}
\subfile{chapters/danh-sach-viet-tat}

% ------------------- PHẦN TÓM TẮT --------------------
\clearpage
\pagenumbering{arabic}          % Chuyển sang đánh số thứ tự: 1, 2, 3...


% ------------------- PHẦN NỘI DUNG CHÍNH ----------
\chapter{GIỚI THIỆU VẤN ĐỀ}
\subfile{chapters/chapter01}    % Gọi nội dung file Chương 1

\chapter{CƠ SỞ LÝ THUYẾT}
\subfile{chapters/chapter02}    % Gọi nội dung file Chương 2

\chapter{PHÂN TÍCH VÀ THIẾT KẾ}
\subfile{chapters/chapter03}    % Gọi nội dung file Chương 3

\chapter{TRIỂN KHAI ỨNG DỤNG}
\subfile{chapters/chapter04}    % Gọi nội dung file Chương 4

\chapter{CHẠY THỬ NGHIỆM VÀ ĐÁNH GIÁ}
\subfile{chapters/chapter05}    % Gọi nội dung file Chương 5

\chapter{TỔNG KẾT}
\subfile{chapters/chapter06}    % Gọi nội dung file Chương 6

% ------------------- PHẦN TÀI LIỆU THAM KHẢO ------
\clearpage
\phantomsection
\addcontentsline{toc}{chapter}{TÀI LIỆU THAM KHẢO}
\bibliography{mybib}            % Trích xuất từ file mybib.bib
\bibliographystyle{plain}       % Định dạng xuất danh sách chuẩn

\end{document}
